\documentclass{article}

\usepackage{amsmath}
\usepackage{amssymb}

\newcommand{\ZZ}{\mathbb{Z}}
\newcommand{\RR}{\mathbb{R}}
\newcommand{\CC}{\mathbb{C}}

\begin{document}

{\bf Worksheet \#13; date: 10/10/2018}

{\bf MATH 55 Discrete Mathematics}

\begin{enumerate}
\item {\em (Rosen 5.3.6b, e)} Determine whether each of these proposed definitions is a valid recursive definition of a function $f$ from the set of nonnegative integers to the set of integers. If $f$ is well defined, find a formula for $f(n)$ when $n$ is a nonnegative integer and prove that your formula is valid.
\begin{enumerate}
\item $f(0) = 1, f(1) = 0, f(2) = 2, f(n) = 2f(n-3)$ for $n \ge 3$
\item $f(0) = 2, f(n) = f(n-1)$ if $n$ is odd and $n \ge 1$ and $f(n) = 2f(n-2)$ if $n \ge 2$
\end{enumerate}

\item {\em (Rosen 5.3.26a, c)} Let $S$ be the subset of the set of ordered pairs of integers defined recursively by

{\em Basis step.} $(0, 0) \in S$

{\em Recursive step.} If $(a, b) \in S$, then $(a + 2, b + 3) \in S$ and $(a + 3, b + 2) \in S$.

\begin{enumerate}
\item List the elements of $S$ produced by the first five replications of the recursive definition.
\item Use structural induction to show that $5 | a + b$ when $(a, b) \in S$.
\end{enumerate}

\item {\em (Challenging)} Show that all Boolean function of $n$ variables can be expressed using just $\wedge$ and $\neg$.

\item {\em (Rosen 6.1.27)} A committee is formed consisting of one representative from each of the $50$ states in the United States, where the representative from a state is either the governor or one of the two senators from that state. How many ways are there to form this committee?

\item {\em (Thank Wikipedia for this example)} California license plate consist of one digit, followed by three letters and 3 digits. However, I, O and Q are only used as the second letter. How many possible plates are there?

\item {\em (Rosen 6.1.47)} In how many ways can a photographer at a wedding arrange six people in a row, including the bride and the groom, if
\begin{enumerate}
\item the bride must be next ot the groom?
\item the bride is not next to the groom?
\item the bride is positioned somewhere to the left of the groom?
\end{enumerate}

\item {\em (Rosen 6.1.65)} How many ways are there to arrange the letters $a$, $b$, $c$, and $d$ such that $a$ is not followed immediately by $b$?

\item {\em (Challenging)} What are the number of derangements of the first $n$ integers, i.e.\ number of arrangements such that $1$ does not come first, $2$ does not come second, $3$ does not come third, etc.?
\end{enumerate}

\end{document}